\documentclass[10pt,a4paper]{article}
\usepackage[margin=25mm]{geometry}
\usepackage{amsmath,amssymb,amsthm,booktabs,array}
\usepackage[T1]{fontenc}
\usepackage{lmodern,microtype,xurl}
\usepackage[hidelinks]{hyperref}
\setlength{\parskip}{4pt}
\setlength{\emergencystretch}{2em}
\renewcommand{\arraystretch}{1.12}
\newtheorem{theorem}{Theorem}
\newtheorem{lemma}[theorem]{Lemma}
\newtheorem{proposition}[theorem]{Proposition}
\theoremstyle{definition}\newtheorem{definition}[theorem]{Definition}
\newcommand{\C}{\mathbb C}
\newcommand{\coef}[1]{[#1]Q}
\title{Global identifiability of a behavioral SIR model\newline
for every fixed real saturation exponent $k\geq2$}
\author{Tatsuhiko Takimoto\\\small Independent researcher}
\date{Version 3.0, 27 October 2026\\Publication display edition, 10 November 2026}
\begin{document}
\maketitle
\begin{abstract}
We prove generic global structural identifiability of the positive real
behavioral SIR model for each fixed, known real exponent $k\geq2$, from
exact observation of the infectious population and the known initial
total population. This proves a specified formulation of Conjecture~2
of Pant, Saucedo and Pogudin: their conjecture sentence does not
explicitly stipulate that the exponent is known or that the domain is
positive, whereas these are assumptions of the theorem here.
For every fixed exponent, an open dense set of true points is
identified against all positive competing points at that same exponent.
For rational $k=p/q$, a root substitution gives a two-state rational
model whose minimal input--output equation has order two and degree
$q$ in its highest derivative. For irrational $k$, the original
supplement's auxiliary-state model has no rational constants beyond
the parameter field. A primitive third-order equation and a finite
Wronskian then identify the rates and hidden states. Both arguments
exclude exceptional exponents throughout $k\geq2$.
The proof is in prose. Separate exact computations check identities
and finite cases. The exponent itself is not inferred, and no claim
of stable estimation from noisy or discrete observations is made.
\end{abstract}

\section{Model, prior work and result}
\label{sec:intro}
Consider
\begin{equation}
S'=-\frac{\beta SI}{(1+\alpha I)^k},\qquad
I'=\frac{\beta SI}{(1+\alpha I)^k}-\gamma I,\qquad
R'=\gamma I,\qquad y=I.
\label{eq:model}
\end{equation}
The exponent $k$ and the initial total population
$N_0=S(0)+I(0)+R(0)>0$ are fixed and known. Parameters
$\alpha,\beta,\gamma$ and initial states are positive. Powers use the
positive real branch. The unknown quantities are the three parameters
and the initial states; the state trajectories then follow by uniqueness
of the analytic initial-value problem. Structural identifiability is
an assertion about exact output germs, not numerical stability.

Pant, Saucedo and Pogudin posed this question in Section~3.6 and
Conjecture~2 of \cite{PSP}. Their Theorem~3.8 gives a two-fold
ambiguity at $k=1$, while Supplement~S9 \cite{PSPsupp} reports global
identifiability for $k=2,3,4,5,10,50$. Its Remark~1 already introduces
\begin{equation}
Z=\frac{(1+\alpha I)^k}{\beta}
\label{eq:S9}
\end{equation}
and uses a computed identifiability result for the transformed model
and the possible zeros of recovery denominators as functions of $k$.
That argument has no integer restriction but specifies neither a
threshold nor all exceptional exponents. The transformation itself
is due to that supplement.

Here the rational and irrational cases are proved separately, with
no exceptional exponent in the stated range. This proves a specified formulation of Conjecture~2,
with each real exponent fixed and known, on the
positive real domain with known $N_0$; the original conjecture sentence
does not explicitly give the first two restrictions. The domain choice
is substantive: for non-integer rational $k=p/q$, the complex cover
identifies $\beta$ only up to $q$ distinct values
(Section~\ref{sec:identification}), whereas the positive branch selects
one. The surrounding
fixed-exponent examples motivate this formulation, but we do not infer
from them a claim about the authors' intended treatment of unknown $k$.
The integer theorem of \cite{V1} is included as a special case on the
positive domain. The rational proof uses the identifiability results
of Ovchinnikov, Pogudin and Thompson \cite{OPT} and the model-ideal
framework of Hong, Ovchinnikov, Pogudin and Yap \cite{HOPY}.
The irrational proof below gives a direct Wronskian argument for the
comparison with every positive competitor.

\begin{samepage}
\begin{definition}
\label{def:id}
Fix $k$ and $N_0>0$. Let $\Omega_{N_0}$ be the open set of positive
$(\alpha,\beta,\gamma,S_0,I_0)$ with $S_0+I_0<N_0$, and put
$R_0=N_0-S_0-I_0$. A quantity is generically globally identifiable if
there is an open dense set $U\subset\Omega_{N_0}$ such that, for each
true point in $U$, every point of $\Omega_{N_0}$ producing the same
analytic output germ at time zero has the same value of that quantity.
The competing points need not belong to $U$.
\end{definition}
\end{samepage}

\begin{theorem}
\label{thm:main}
For every fixed, known real $k\geq2$ and every fixed, known $N_0>0$,
model~\eqref{eq:model} is generically globally identifiable in the
sense of Definition~\ref{def:id}. Thus $\alpha,\beta,\gamma,S_0,I_0,R_0$
are determined by $I(t)$ and $N_0$ at every true point of an open dense
set $U_{k,N_0}\subset\Omega_{N_0}$, against all positive competitors
at the same exponent and population. The set may depend on $k$ and $N_0$.
\end{theorem}

\paragraph{Present assurance.}
The theorem has a prose proof and is not formalized in a proof assistant
or reviewed by human experts. Separate AI proof checks and separately
written exact computations supported the rational and irrational
arguments, including their all-competitor quantifiers. The finite
calculations in Section~\ref{sec:checks} check algebraic identities and
specified examples; they are not the proof for all real exponents.
This extract contains the article TeX and PDF. The companion
\texttt{README.md} specifies its scope and distinguishes the original
supplement and verification record from those of the minimal-jet packet.
\section{A rational cover and its constants}
\label{sec:cover}
In this and the next three sections, write $k=p/q\geq2$ in lowest terms, with $p,q$ positive integers. Set $K=\C(\alpha,\beta,\gamma)$ and let $L=K(S,z)$, with $S,z$
algebraically independent. Define the output and derivation by
\begin{equation}
y=\frac{z^q-1}{\alpha},\qquad
DS=-\frac{\beta Sy}{z^p},\qquad
Dz=\frac{\alpha}{q}z^{1-q}
       \left(\frac{\beta Sy}{z^p}-\gamma y\right),\qquad D|_K=0.
\label{eq:cover}
\end{equation}
This is a rational ODE with rational output. A common denominator is
$\alpha qz^{p+q-1}$. On the positive real domain the changes of initial
coordinates
\[
z=(1+\alpha I)^{1/q}>0,\qquad I=(z^q-1)/\alpha
\]
are analytic inverses, and $z^p=(1+\alpha I)^k$. Differentiation shows
that they carry solutions in both directions. In particular, the cover
does not add a free initial condition to the positive real model.
The equation for $R$ can be removed, since $R=N_0-S-I$ and
$D(N_0-S-I)=\gamma I$.

\begin{lemma}
\label{lem:constants}
The field of constants of $(L,D)$ is exactly $K$.
\end{lemma}
\begin{proof}
The derivation $E=(z^p/y)D$ has the same constants as $D$, and
\begin{equation}
ES=-\beta S,\qquad
Ez=\frac{\alpha\beta}{q}Sz^{1-q}
       -\frac{\alpha\gamma}{q}z^{p-q+1}.
\label{eq:E}
\end{equation}
Equivalently,
\[
E=-\beta S\partial_S+
\left(\frac{\alpha\beta}{q}Sz^{1-q}
-\frac{\alpha\gamma}{q}z^{p-q+1}\right)\partial_z.
\]
For a nonzero $F\in K(z)(S)$ with $EF=0$, use the embedding in the
formal Laurent series field $K(z)((S))$ and write
$F=\sum_{n\geq m}f_n(z)S^n$, $f_m\ne0$. The derivation does not lower
$S$-order, so it extends continuously to $K(z)((S))$; its term
containing $S\partial_z$ raises that order by one.
Thus the coefficient of $S^m$ in $EF$ gives
\begin{equation}
-\beta m f_m-\frac{\alpha\gamma}{q}z^{p-q+1}f'_m=0.
\label{eq:leading}
\end{equation}
If $m\ne0$, then
\[
\frac{f'_m}{f_m}=-\frac{q\beta m}{\alpha\gamma}z^{-(p-q+1)}.
\]
The logarithmic derivative of a nonzero rational function has at most
a simple pole at $z=0$: writing $f_m=z^r h$ with $h$ of valuation zero
gives $f'_m/f_m=r/z+h'/h$. The right side has pole order
$p-q+1\geq q+1\geq2$, a contradiction. Hence $m=0$ and
\eqref{eq:leading} gives $f_0\in K$. If $F-f_0$ were nonzero, it would
again be a constant with positive $S$-order, giving the same
contradiction. Therefore $F=f_0\in K$.
\end{proof}

\section{Differential order two and highest-derivative degree \texorpdfstring{$q$}{q}}
\label{sec:io}
Put $u=1+\alpha y$, $v=Dy$ and $w=D^2y$, recalling $k=p/q$. Define
\begin{equation}
A=yu w-\bigl(1-(k-1)\alpha y\bigr)v^2+k\alpha\gamma y^2v.
\label{eq:A}
\end{equation}
The symbol $A$ denotes this polynomial in output derivatives; it is
not the auxiliary state $Z$ in~\eqref{eq:S9}.
From the output equation and conservation we have
\begin{equation}
S=\frac{(v+\gamma y)z^p}{\beta y},\qquad DS=-(v+\gamma y).
\label{eq:S}
\end{equation}
Differentiate the first expression and use $D(z^q)=\alpha v$ and the
second expression. Multiplication by $\beta y^2$ gives
\[
(w+\gamma v)z^p y+k\alpha vy(v+\gamma y)z^{p-q}
-(v+\gamma y)z^p v=-\beta y^2(v+\gamma y).
\]
Using $u-k\alpha y=1-(k-1)\alpha y$ and collecting terms gives
\begin{equation}
z^{p-q}A+\beta y^2(v+\gamma y)=0.
\label{eq:unpowered}
\end{equation}
Raising this equality to the $q$th power yields
\begin{equation}
Q(y,v,w)=u^{p-q}A^q-(-\beta)^q y^{2q}(v+\gamma y)^q=0.
\label{eq:Q}
\end{equation}
Keeping the factor $(-\beta)^q$ makes the sign correct for odd as well
as even denominators.

\begin{lemma}
\label{lem:minimal}
The elements $y,v$ are algebraically independent over $K$. The
polynomial $Q$ is irreducible in $K[y,v,w]$ and has minimal degree
$q$ in $w$ over $K(y,v)$.
\end{lemma}
\begin{proof}
Equation~\eqref{eq:S} gives $K(S,z)=K(v,z)$, so $v,z$ are independent;
then $y,v$ are independent as well. Write $F=K(y,v)$. The polynomial
$T^q-u$ is Eisenstein at $u$ in $K(v)[u][T]$ (and is linear when
$q=1$), hence $[F(z):F]=q$. Its conjugates are $z\mapsto\zeta z$ for
$\zeta^q=1$, since $\C$ contains the roots of unity.
Equation~\eqref{eq:unpowered} writes $w$ as an element of $F$ plus a
nonzero element of $F$ times $z^{-(p-q)}$. Because
$\gcd(p-q,q)=1$, its $q$ conjugates are distinct. Therefore $w$ has
minimal degree $q$ and $Q$ is irreducible in $F[w]$.

The leading coefficient of $Q$ as a polynomial in $w$ is $y^q u^p$.
Its content in $K[y,v][w]$ can thus have only the factors $y$ or $u$.
However,
\begin{equation}
Q(0,v,w)=(-1)^qv^{2q},\qquad
Q(-1/\alpha,v,w)=-(-\beta)^q\alpha^{-2q}(v-\gamma/\alpha)^q
\label{eq:content}
\end{equation}
are nonzero; the second equality uses $p-q>0$.
Thus $Q$ is primitive, and Gauss's lemma proves irreducibility in
$K[y,v,w]$.
\end{proof}

\begin{proposition}
\label{prop:characteristic}
For the ranking $y<y'<y''<\cdots$, the set $\{Q\}$ is a monic
characteristic presentation of the output ideal of~\eqref{eq:cover}.
\end{proposition}
\begin{proof}
The output ideal $J$ is the kernel of $K\{y\}\to L$ and is prime.
This agrees with the output part of the rational-model ideal obtained
by saturating by the common denominator: after localization, all
state derivatives reduce to rational expressions in the independent
states $S,z$. This is also the model-ideal setting of
\cite[Lemma~3.2]{HOPY}.

The leader of $Q$ is $y''$, with degree $q$. Lemma~\ref{lem:minimal}
excludes any nonzero element of $J$ of lower rank. An autoreduced set
in one differential indeterminate has at most one element; hence
$\{Q\}$ is a characteristic set. Its initial $y^q u^p$ belongs to
$K[y,v]$, the ring of non-leading variables. Separately, $Q$ itself
has no nonconstant factor in $K[y,v]$, by primitivity.
Finally $\coef{y^q w^q}=1$. These are the conditions in
\cite[Definition~5.4]{OPT}. Here ``monic'' means that at least one
monomial coefficient is one, not that the initial is one.
\end{proof}

\section{Coefficients and the integer special case}
\label{sec:coefficients}
\begin{proposition}
\label{prop:coefficients}
For $q\geq2$, the field generated over $\C$ by the coefficients of $Q$
is $\C(\alpha,\gamma,\beta^q)$. For $q=1$ it is $K$.
\end{proposition}
\begin{proof}
When $q\geq2$, direct expansion gives
\begin{equation}
\begin{aligned}
\coef{y^qw^q}&=1,&\quad \coef{y^{q+1}w^q}&=p\alpha,\\
\coef{y^{q+1}vw^{q-1}}&=p\alpha\gamma,&
\coef{y^{3q}}&=(-1)^{q+1}\beta^q\gamma^q.
\end{aligned}
\label{eq:coeff}
\end{equation}
For the third coefficient, choose the $w$ term $q-1$ times and the
$k\alpha\gamma y^2v$ term once from $A^q$, taking the least $y$ power
in the remaining factors. Terms containing $v^2$ cannot contribute.
The second summand of~\eqref{eq:Q} has no $w$, and $q-1>0$.
Thus
\begin{equation}
\alpha=\frac{\coef{y^{q+1}w^q}}p,\quad
\gamma=\frac{\coef{y^{q+1}vw^{q-1}}}{p\alpha},\quad
\beta^q=\frac{(-1)^{q+1}\coef{y^{3q}}}{\gamma^q}.
\label{eq:inverse}
\end{equation}
All coefficients lie in this field, so equality follows.

If $q=1$, $k=p\geq2$ and $Q=u^{p-1}A+\beta y^2(v+\gamma y)$,
which is the integer-exponent polynomial of \cite{V1}. In this case
the $\beta$ term contributes to $\coef{y^2v}$, so one instead uses
\begin{equation}
\coef{yw}=1,\quad \coef{y^2w}=p\alpha,\quad
\coef{y^3v}=p(p-1)\alpha^2\gamma,\quad
\coef{y^2v}=\beta+p\alpha\gamma.
\label{eq:integer}
\end{equation}
These recover $\alpha,\gamma,\beta$ successively. The integer factors
in the denominators of~\eqref{eq:inverse} and~\eqref{eq:integer} never
vanish in their respective ranges. The other divisions are legitimate
at generic points and on the positive parameter domain.
\end{proof}

\section{One experiment and the positive branch}
\label{sec:identification}
By Proposition~\ref{prop:characteristic} and
\cite[Corollary~5.7]{OPT}, the coefficient field is the field of
input--output identifiable parameter functions. Lemma~\ref{lem:constants}
verifies the hypothesis of \cite[Theorem~4.7]{OPT}, so it is also the
field of functions identifiable from one experiment. In particular,
$\alpha,\gamma,\beta^q$ are generically globally identifiable on the
complex rational cover.

The quantifiers here are those of \cite[Notation~2.1 and
Definition~2.2]{OPT}, which extend the definition cited there from
\cite{HOPY}. For this autonomous cover, let
$\Omega_{\mathrm c}=\{(\alpha,\beta,\gamma,S_0,z_0):
\alpha qz_0^{p+q-1}\ne0\}$. For each identifiable function $h$,
the true point lies in a nonempty Zariski open set
$\Theta\subset\Omega_{\mathrm c}$, while every competing point in
$\Omega_{\mathrm c}$ where $h$ is defined and the output germ is the same
must give the same value of $h$. Competitors are not restricted to
$\Theta$. For $h=\alpha,\gamma,\beta^q$ these functions are defined
everywhere, and we intersect their three Zariski open sets. This is
the quantifier structure specialized to the positive branch in
Definition~\ref{def:id}.

For $q>1$, $\beta$ itself is not globally identifiable on that cover.
For each $\zeta^q=1$, the transformation
\begin{equation}
(z,\beta,S,\alpha,\gamma)\longmapsto
(\zeta z,\zeta^p\beta,S,\alpha,\gamma)
\label{eq:symmetry}
\end{equation}
preserves the output and $DS$ and multiplies $Dz$ by $\zeta$.
It therefore takes solutions to solutions. Coprimality gives $q$
distinct generic values of $\beta$, and identifiability of $\beta^q$
allows no more than $q$.

\begin{proof}[Rational case of Theorem~\ref{thm:main}]
Take a true point in the Zariski open set $\Theta$ just described for the complex-cover
identifiability statement, and any positive real competing point with
the same output. Both lift by~\eqref{eq:cover}. The global statement
compares the true point with all admissible competitors, so
$\alpha,\gamma,\beta^q$ agree. Their positive roots
$z=(1+\alpha y)^{1/q}$ agree as well. The relation before taking
powers gives
\begin{equation}
\beta=-\frac{z^{p-q}A}{y^2(v+\gamma y)}.
\label{eq:beta}
\end{equation}
This recovers the same $\beta$ for both solutions. Its denominator is
nonzero when $\beta Sy z\ne0$, since $v+\gamma y=\beta Sy/z^p$.
For positive $\beta$ one could alternatively take the positive root
of $\beta^q$; equation~\eqref{eq:beta} displays the role of the branch.

Combining~\eqref{eq:S} with~\eqref{eq:unpowered} gives
\begin{equation}
S=-\frac{yu(v+\gamma y)^2}{A},\qquad I=y,\qquad R=N_0-S-I.
\label{eq:states}
\end{equation}
The right sides coincide. In fact, along every positive solution,
\eqref{eq:unpowered} gives
$z^{p-q}A=-\beta^2Sy^3/z^p\ne0$. Thus both solutions have the same
nonzero denominator, proving equality of all initial states.
The rational expression for $S$ also gives its identifiability on the
cover by \cite[Proposition~3.1]{OPT}.

For each fixed $p,q$, intersect $\Theta$ with the
nonvanishing sets of the displayed denominators. This is a nonempty
Zariski open set, by the cited definition and finite intersection.
The positive real cover with $S+I<N_0$ is a nonempty
Euclidean open set in real coordinates. A nonzero complex polynomial
cannot vanish on that whole open set, so the intersection is open and
dense there. The analytic coordinate change in
Section~\ref{sec:cover} transfers it to $\Omega_{N_0}$. This restricts
the true point only; it does not restrict the competing points to a
generic subset.
\end{proof}

For $q=1$, no root choice is needed. The same coefficient and state
formulas recover the complex generic integer result of \cite{V1}.
The new real-branch theorem therefore includes its positive real
restriction and does not assert a contrary ambiguity for integers.


\section{The irrational auxiliary-state field}
\label{sec:irrational}
Now fix an irrational real $k\geq2$, and use the state $Z$ of
\eqref{eq:S9}. In this section and the next two, redefine
$K=\C(\alpha,\gamma)$, in contrast to $K=\C(\alpha,\beta,\gamma)$
in the rational case: $\beta$ is encoded by $Z(0)$, not included as an
independent parameter of this field. Set $L=K(S,u,Z)$ with three
algebraically independent states. The output
and derivation are
\begin{equation}
\begin{aligned}
y&=(u-1)/\alpha, & DS&=-Sy/Z,\\
Du&=\alpha y(S/Z-\gamma), &
DZ&=k\alpha y(S-\gamma Z)/u.
\end{aligned}
\label{eq:irrmodel}
\end{equation}
The parameter $\beta$ does not occur in the vector field: it is encoded
by the initial value $Z(0)=u(0)^k/\beta$. This is the model of S9,
with a changed name for its auxiliary state.

\begin{lemma}
\label{lem:irrconstants}
For each fixed irrational $k\geq2$, the field of constants is
$L^D:=\{H\in L:DH=0\}=K$.
\end{lemma}
\begin{proof}
Rescale by the nonzero field element $Z/y$ and set $B=Z/u$. The
derivation $E=(Z/y)D$ has the same constants and is
\begin{equation}
E=-S\partial_S-\alpha\gamma B X+
 \alpha S\left(\partial_u+\frac{(k-1)B}{u}\partial_B\right),
\qquad X=u\partial_u+(k-1)B\partial_B.
\label{eq:irrE}
\end{equation}
It does not lower $S$-order and extends continuously to
$K(u,B)((S))$. For a nonzero rational $H$ with $EH=0$, write
$H=\sum_{n\geq m}h_n(u,B)S^n$, $h_m\ne0$. Its leading coefficient
satisfies
\[
E_0h_m=mh_m,\qquad E_0=-\alpha\gamma B X.
\]
Expand $h_m=\sum_{j\geq n}g_j(u)B^j$ in $K(u)((B))$. Since
\[
E_0(g_jB^j)=-\alpha\gamma\bigl(ug'_j+(k-1)jg_j\bigr)B^{j+1},
\]
$E_0$ raises the lowest $B$-order by at least one. The equality is
impossible for $m\ne0$. Thus $m=0$ and $Xh_0=0$.

Write $h_0=P_0/Q_0$ with coprime polynomials in $K[u,B]$.
The equality $Q_0XP_0=P_0XQ_0$ implies that $P_0$ divides $XP_0$
and $Q_0$ divides $XQ_0$. Since $X$ does not increase total degree,
$XP_0=cP_0$ and $XQ_0=cQ_0$ for the same $c\in K$.
The weights $i+(k-1)j$ of distinct monomials $u^iB^j$ are distinct:
a coincidence with different $j$ would make $k$ rational, and equal
$j$ forces equal $i$. Each eigenpolynomial is therefore a monomial.
The common eigenvalue gives the same exponents, so $h_0\in K$.
If $H-h_0$ were nonzero, it would be an $E$-constant of positive
$S$-order, contradicting the preceding argument. Hence $H\in K$.
\end{proof}

The irrationality is used only in the monomial weights. For rational
$k=p/q$, the three-state field instead has the nonconstant rational
first integral $u^p/Z^q$. This does not contradict
Lemma~\ref{lem:constants}: on the two-state cover that expression is
$\beta^q$, in the parameter field. For every real $k$, the positive
model also has the analytic first integral $Z/u^k$. In the irrational
case it is not rational. Indeed, a putative polynomial relation for
$u^k$ over $\C(u)$ gives a finite linear combination of distinct powers
$u^{i+kj}$; after $u=e^s$, a Vandermonde argument for exponentials
forces every coefficient to vanish. The lemma concerns rational
constants, not all analytic first integrals.

\section{Third-order equation and coefficient recovery}
\label{sec:irrIO}
Use the jet variables $v=Dy$, $w=D^2y$ and $x=D^3y$.
Retain the polynomial $A$ from~\eqref{eq:A}, now with the
fixed irrational value of $k$, and put $r=v+\gamma y$ and $u=1+\alpha y$.
The original equations give $S=Zr/y$, $DS=-r$ and $DZ/Z=k\alpha v/u$.
Differentiating the first identity and multiplying by $y^2u$ gives
\[
Z\{k\alpha vyr+uy(w+\gamma v)-urv\}=-y^2ur.
\]
The expression in braces is $A$, so
\begin{equation}
ZA=-y^2ur,\qquad Z=-\frac{y^2ur}{A},\qquad
S=-\frac{yur^2}{A},\qquad K(S,u,Z)=K(y,v,w).
\label{eq:irrrecover}
\end{equation}
In particular $y,v,w$ are algebraically independent. Logarithmic
differentiation of $Z$ yields
\[
\frac{DA}{A}=\frac{2v}{y}+\frac{w+\gamma v}{r}
 -(k-1)\frac{\alpha v}{u}.
\]
Define
\begin{equation}
\begin{split}
P={}&uyr\,(v\partial_yA+w\partial_vA+x\partial_wA)\\
 &-A\{2uvr+uy(w+\gamma v)-(k-1)\alpha vyr\}.
\end{split}
\label{eq:irrP}
\end{equation}
It vanishes in $L$ and is linear in $x$, with
\begin{equation}
\partial_xP=y^2u^2r.
\label{eq:irrinitial}
\end{equation}

\begin{lemma}
\label{lem:irrminimal}
For each fixed $k\ne0,1$, the polynomial $P$ is primitive and
irreducible in $K[y,v,w,x]$. Its differential order is minimally three
and its degree in the highest derivative is minimally one. Its
coefficients generate $K=\C(\alpha,\gamma)$.
\end{lemma}
\begin{proof}
The content in $K[y,v,w][x]$ can have only the factors $y,u,r$.
They are excluded by
\begin{equation}
\begin{aligned}
P|_{y=0}&=2v^4,\\
P|_{y=-1/\alpha}&=k(k-1)v^2(v-\gamma/\alpha)^2,\\
P|_{v=-\gamma y}&=-y^2u^2(w-\gamma^2y)^2.
\end{aligned}
\label{eq:irrcontent}
\end{equation}
These are nonzero for $k\ne0,1$. Primitivity and linearity in $x$
give irreducibility by Gauss's lemma. The field equality in
\eqref{eq:irrrecover} proves that no lower-order relation exists;
linearity gives the minimal positive degree. Moreover the kernel of
$K[y,v,w,x]\to L$ is exactly $(P)$: over $K(y,v,w)$ it is generated
by the linear equation, and primitivity contracts this ideal to $(P)$.
The highest-derivative terms give
\begin{equation}
[y^2vx]P=1,\qquad [y^3vx]P=2\alpha,\qquad [y^3x]P=\gamma.
\label{eq:irrcoeff}
\end{equation}
All coefficients belong to $K$, and these recover its generators.
\end{proof}

For the differential ranking by order, $\{P\}$ is also a monic
characteristic presentation: it has minimal rank in the prime output
ideal, its initial is in the non-leading-variable ring, it has no
nonconstant factor there, and a monomial coefficient is one.
These are the same conditions as in Proposition~\ref{prop:characteristic}.
One may therefore use \cite[Corollary~5.7 and Theorem~4.7]{OPT} with
Lemma~\ref{lem:irrconstants}. The following direct argument spells out
why only the true point needs to be generic.

\section{Wronskians and all positive competitors}
\label{sec:wronskian}
We use the characteristic-zero Wronskian criterion: elements of a
differential field are linearly independent over its constants if
and only if their Wronskian is nonzero. For completeness, divide
columns by a nonzero first element. Elementary row operations give
$W(f_1,\ldots,f_n)=f_1^n W(1,f_2/f_1,\ldots,f_n/f_1)$; expansion
along the first column reduces to the Wronskian of
$D(f_2/f_1),\ldots,D(f_n/f_1)$. If it vanishes, induction gives a
nontrivial constant linear combination of those derivatives equal
to zero. Its antiderivative is a constant, giving a relation among
the original elements. The converse follows by differentiating a
constant linear relation. The one-element case starts the induction.

For the fixed irrational $k$, let $\mathcal M$ be the union, as
$\alpha,\gamma$ vary, of the output-monomial supports of $P$.
Expansion of~\eqref{eq:irrP} gives $|\mathcal M|=25$ for irrational
$k$, so the Wronskian below has $n=24$ columns. Some coefficients
vanish at rational exponents (for example, $k=2$), but this section
fixes an irrational exponent. All monomials have total degree at
most six.
The polynomial $P$ has total degree exactly six, as seen from the
term $\alpha^2\gamma y^5x$ in~\eqref{eq:irrinitial}.
A $K$-linear relation among the evaluations of $\mathcal M$ in $L$
is a polynomial in the kernel $(P)$ of Lemma~\ref{lem:irrminimal}.
Its degree is at most six, so it is a scalar multiple of $P$.
There is exactly one relation up to scalar.

Remove $m_0=y^2vx$, whose coefficient in $P$ is one. The remaining
monomials $m_1,\ldots,m_n$ are $K$-linearly independent in $L$.
By Lemma~\ref{lem:irrconstants} and the Wronskian criterion,
\begin{equation}
W=\det\bigl(D^{i-1}m_j\bigr)_{1\leq i,j\leq n}\ne0
\quad\hbox{in }K(S,u,Z).
\label{eq:W}
\end{equation}
This proves nonvanishing as a rational function; a numerical
Wronskian calculation is not needed.

Restrict the true initial point so that $W$ is defined and nonzero.
Include the common denominator $\alpha uZ$ and the denominators used
to write $W$; their nonvanishing defines a nonempty Zariski open set
in the transformed parameter-state space. Any positive competitor
with the same output obeys its own polynomial $P_*$, with parameters
$\alpha_*,\gamma_*$ but the same fixed $k$. Its support is contained
in $\mathcal M$, and its coefficient of $m_0$ is one. Thus $P_*-P$
is a constant linear combination of $m_1,\ldots,m_n$ along the true
output. Differentiate this zero relation $0,\ldots,n-1$ times and
evaluate at time zero. The coefficient differences lie in the kernel
of the invertible Wronskian matrix, so all are zero.
Equation~\eqref{eq:irrcoeff} gives $\alpha_* =\alpha$ and
$\gamma_* =\gamma$. No condition on the competitor's own Wronskian
was imposed.

\begin{proof}[Irrational case and completion of Theorem~\ref{thm:main}]
The positive initial-coordinate map
\begin{equation}
(\alpha,\beta,\gamma,S_0,I_0)\longmapsto
(\alpha,\gamma,S_0,u_0=1+\alpha I_0,Z_0=u_0^k/\beta)
\label{eq:irrmap}
\end{equation}
is analytically invertible, with inverse
$I_0=(u_0-1)/\alpha$ and $\beta=u_0^k/Z_0$.
It carries positive solutions of~\eqref{eq:model} to
\eqref{eq:irrmodel} and back: differentiating $u^k/Z$ along the latter
system gives zero. The image domain is the nonempty real open set
$\alpha,\gamma,S_0,Z_0>0$, $u_0>1$ and
$S_0+(u_0-1)/\alpha<N_0$. A nonzero complex polynomial cannot vanish
on that entire open set. Hence the Wronskian's nonvanishing set is
open dense there, and its inverse image under~\eqref{eq:irrmap} is
open dense in $\Omega_{N_0}$.

For a true point there and any positive competitor, the preceding
argument gives the same $\alpha,\gamma$. Equation~\eqref{eq:irrrecover}
then gives the same $S,Z,u$ from the common output jets. It is defined
at every positive competitor, since
$y>0$, $u>1$, $r=Sy/Z>0$ and $A=-y^2ur/Z<0$.
Finally
\[
\beta=\frac{u_0^k}{Z_0},\qquad I_0=y(0),\qquad
R_0=N_0-S_0-I_0
\]
agree. Together with the rational case, this proves the theorem for
every real $k\geq2$. No limiting argument from rational exponents is
used.
\end{proof}
\section{Exact finite checks}
\label{sec:checks}
For rational exponents, two sets of computations supply 22 case records at 19 distinct exponents,
18 of them in the theorem's range. The producer side consists of a
SymPy program for symbolic rational-function identities and a separate
standard-library program for sparse polynomials and finite-field
Taylor matrices. It covers
\[
\frac32,\frac52,\frac72,\frac92,\frac{11}2,\frac{13}2,
\frac{15}2,\frac{17}2,\frac{19}2,\frac73,\frac83,\frac{11}4,2,3.
\]
For each case, it checks~\eqref{eq:unpowered}--\eqref{eq:Q}, the
coefficient recovery, primitivity substitutions,
\eqref{eq:E}, \eqref{eq:symmetry} and~\eqref{eq:states}.
The second set uses separately written sparse Laurent polynomials over
$\mathbb Q$, with parameters left symbolic. It covers
\[
\frac{13}6,\frac{17}5,\frac{23}7,\frac94,\frac{31}8,2,3,\frac52
\]
and checks the same algebraic identities and selected coefficients.
The three overlapping exponents are regression cases. The value
$3/2$ is a comparison below the theorem's range, not an instance of
Theorem~\ref{thm:main}.

Only the first 14 cases have finite-field observation matrices.
At $(\alpha,\beta,\gamma,S,z)=(2,3,5,7,2)$, Taylor series are computed
exactly modulo $1000003$. If $m$ is the number of output monomials
in $Q$, their series are retained through degree $m+5$ (state series
through $m+7$). All denominators are invertible modulo this prime.
The column for $y^q w^q$, whose coefficient is one, is removed;
the remaining $m-1$ columns have full rank in every case. For
$k=5/2$ the rank is $31$, and for $k=11/4$ it is $161$.
The matrices and nonzero minors belong to the original companion
supplement, outside this article-only extract (see its \texttt{README.md}).

A nonzero minor at this point modulo the prime certifies that the
corresponding rational-function minor is not identically zero.
On its generic nonvanishing set, any two competing coefficient
vectors lie in the kernel of the same observation matrix and share
the normalized coefficient one. Their difference has that coefficient
zero; independence of the other columns makes the difference zero.
This certifies generic coefficient uniqueness for these fixed
exponents. It is distinct from extrapolating a finite list to all
rational $k$; that conclusion rests on the preceding proof.


For irrational exponents, two producer programs evaluate the
three-state identities using sparse Laurent arithmetic over
$\mathbb Q$ and SymPy, respectively. Their full expansions of $P$
agree in all 36 terms (25 output monomials). They check the original
ODE substitution, both recovery identities, the three content
substitutions, and the coefficients in~\eqref{eq:irrcoeff}.
The CAS path specializes symbolically at $\sqrt5$, $\pi$ and
$1+\sqrt2$, without replacing these values by approximate rationals.
The constants-field transformations were also checked on 27
Laurent monomials, totaling 90 expanded coefficient terms, with
81 weight pairs $(i,j)\in\{-4,\ldots,4\}^2$ and rational controls;
the all-rational-function conclusion remains the prose lemma.

A separately written standard-library implementation derives state
jets by formal Taylor integration of~\eqref{eq:irrmodel}, agrees on
all 36 polynomial terms, and works in $\mathbb Q(\sqrt7)$ at
$(\alpha,\gamma,S_0,I_0,Z_0)=(2,3,5,2,4)$, $(2,3,12,2,4)$ and
$(1,1,7,3,2)$. It checks the equation and both recovery identities
through Taylor order eight. The middle case has $y'(0)=0$ and
$y''(0)=-3$, so it includes a local infection maximum. With $N_0=20$
all three cases have positive $R_0$. The original companion supplement,
outside this extract, retains a rejected missing-factor formula as a
negative control: at the
maximum its initial residual vanishes but the next Taylor coefficient
does not. This shows why checking a single value is insufficient.
Neither these examples nor the 27 monomial tests are a computation
of the general Wronskian in~\eqref{eq:W}.

\section{Scope of the conclusion}
\label{sec:scope}
The theorem covers all real exponents $k\geq2$, each fixed and known,
on the positive domain. It neither compares different unknown
exponents nor asserts a common generic set for all $k$. The generic
set in the original coordinates is open dense; Zariski openness in
those coordinates is not claimed. In the irrational case the initial
coordinate change uses $u^k$ and is analytic rather than algebraic.
Unknown $N_0$, identifiability at every initial point, and stable
recovery from finite noisy observations remain outside the result.

The rational proof in fact only requires $p>q$, hence also applies
to fixed rational $1<k<2$: the logarithmic-derivative pole has order
$p-q+1\geq2$, the content calculation uses $p-q>0$, and the recovery
coefficients remain valid; $q=1$ then already means $p\geq2$.
The irrational argument likewise uses only irrationality and
$k\ne0,1$; it remains valid for every irrational real $k$, with the
exclusions automatic for irrational $k$. These observations are scope
notes, not extensions of the main theorem's stated range $k\geq2$.

The auxiliary-state construction and large-exponent argument of S9
are prior work. The contribution here is the explicit treatment of
the constants, minimal equations, exceptional exponents and positive
competitors across the stated range. A bounded literature comparison
found no matching all-real-exponent theorem in the sources examined;
this does not establish priority or exclude other formulations known
to specialists.
\section*{AI-use statement}
AI systems assisted with proof development, separate checks of the
proof and algebra, literature organization, computation and manuscript
drafting. The author is responsible for the mathematical claims and
their stated scope; AI systems are not authors.

{\raggedright
\begin{thebibliography}{9}
\bibitem{PSP}
B. Pant, O. Saucedo and G. Pogudin,
``Uncovering identifiability of epidemiological models: basic
reproduction number and complementary data streams,''
medRxiv, version 2 (25 January 2026),
\href{https://doi.org/10.64898/2026.01.16.26344284}{doi:10.64898/2026.01.16.26344284}.
\bibitem{PSPsupp}
B. Pant, O. Saucedo and G. Pogudin,
Supplementary material to \cite{PSP}, version 2, Sections S8 and S9,
\url{https://www.medrxiv.org/content/medrxiv/early/2026/01/25/2026.01.16.26344284/DC1/embed/media-1.pdf}.
\bibitem{OPT}
A. Ovchinnikov, G. Pogudin and P. Thompson,
``Parameter identifiability and input--output equations,''
\emph{Appl. Algebra Engrg. Comm. Comput.} \textbf{34}(2), 165--182
(2023; published online 2021),
\href{https://doi.org/10.1007/s00200-021-00486-8}{doi:10.1007/s00200-021-00486-8};
\href{https://arxiv.org/abs/2007.14787v2}{arXiv:2007.14787v2}.
The numbering used here is that of the cited arXiv version.
\bibitem{HOPY}
H. Hong, A. Ovchinnikov, G. Pogudin and C. Yap,
``Global identifiability of differential models,''
\emph{Comm. Pure Appl. Math.} \textbf{73}(9), 1831--1879 (2020),
\href{https://doi.org/10.1002/cpa.21921}{doi:10.1002/cpa.21921};
\href{https://arxiv.org/abs/1801.08112v5}{arXiv:1801.08112v5}.
The lemma number refers to the cited arXiv version.
\bibitem{V1}
T. Takimoto,
``Global identifiability of a behavioral SIR model for every integer
saturation exponent $k\geq2$,'' article version 1.1, release 1.0.0
(22 October 2026),
\href{https://doi.org/10.5281/zenodo.23553807}{doi:10.5281/zenodo.23553807}.
\end{thebibliography}
}
\end{document}
